\section{Recitation 7: Recursion}
\label{sec:rec7}

This recitation asks two closely related questions: how to build a \emph{circular} data structure, and how to let something refer to \emph{itself} in its own definition. A plain \code{let} answers neither, for the reason of \Cref{sec:recproblem}: the right-hand side is evaluated before the name is bound. The recitation shows two ways around this. The first ties the knot with boxes, as in \Cref{sec:cycles}. The second uses \emph{fixed-point combinators}, which give recursion with no mutation at all. A challenge at the end uses one of them to implement \code{rec} in the interpreter. The code is plain \code{\#lang racket}, with tests written with \code{check-expect} from \code{test-engine/racket-tests}, except for the challenge, which is \code{plai-typed}.

\subsection{Tying the Knot with Boxes}

\code{self-box} is a box that contains itself, and \code{two-cycle} is two boxes that point to each other, $\code{b1} \to \code{b2} \to \code{b1}$:
\begin{lstlisting}
(define self-box
  (let ([b (box 'dummy)])
    (begin
      (set-box! b b)
      b)))
(define two-cycle
  (let* ([b1 (box 'dummy)]
         [b2 (box b1)])
    (begin
      (set-box! b1 b2)
      b1)))

(check-expect (eq? (unbox self-box) self-box) #t)
(check-expect (eq? (unbox (unbox two-cycle)) two-cycle) #t)
(check-expect (eq? (unbox two-cycle) two-cycle) #f)
\end{lstlisting}
The tests use \code{eq?}, which asks whether two values are the \emph{same} object in memory. Unboxing \code{self-box} once gets back to it; \code{two-cycle} needs two steps, and one step gets to the other box. Racket prints the two values as \code{\#0='\#\&\#0\#} and \code{\#0='\#\&\#\&\#0\#} (\Cref{sec:cycles}).

\begin{exercise}[Recursive functions with boxes]
Implement factorial as \code{fact-box} by putting the function in a box, without letting \code{fact-box} call itself by name.
\end{exercise}
\begin{solution}
This is the factorial of \Cref{sec:cycles}, with the call \code{((unbox fact) n)} at the end:
\begin{lstlisting}
(define (fact-box n)
  (let* ([fact (box 'dummy)]
         [fact-fun (lambda (k)
                     (if (equal? k 0)
                         1
                         (* k ((unbox fact) (- k 1)))))])
    (begin
      (set-box! fact fact-fun)
      ((unbox fact) n))))

(check-expect (fact-box 0) 1)
(check-expect (fact-box 5) 120)
\end{lstlisting}
\end{solution}

\begin{pitfall}[\code{(unbox fact)} is the function, not the call]
\code{(* k (unbox fact) (- k 1))} is also syntactically fine, but it multiplies three things: \code{k}, the function itself, and \code{k - 1}. It fails with \code{*: contract violation, expected: number?, given: \#<procedure:fact-fun>}. The test \code{(fact-box 0)} still passes, because the base case never reaches the multiplication, so only \code{(fact-box 5)} shows the bug. The call needs its own parentheses: \code{((unbox fact) (- k 1))}.
\end{pitfall}

\subsubsection{A loop with no recursive definition} What does this program do?
\begin{lstlisting}
(let ([nonterminating (box 'dummy)])
  (begin
    (set-box! nonterminating (lambda (x) ((unbox nonterminating) x)))
    ((unbox nonterminating) 5)))
\end{lstlisting}
It runs forever, although nothing in it is defined in terms of itself. The loop is in the store, not in the program text. The \code{lambda} reaches the box through its environment, and after \code{set-box!} the box holds that same \code{lambda}. So each call unboxes the function and calls it again. \code{(define (loop x) (loop x))} loops for the same reason: a recursive \code{define} also ties a knot, and Racket ties it for us (\Cref{sec:recboxes}). The call is in tail position, so Racket does not grow the stack; memory use stays flat and the program simply never stops.

\begin{exercise}[Circular linked lists]
Write \code{circ-list}, which takes $n \ge 1$ and returns a circular list of $n$ boxes. Each box holds a pair \code{(cons value next)}. The head holds $n$, the next box $n - 1$, and so on down to 1, and the box holding 1 points back to the head:
\[
\text{head } (n) \to (n-1) \to \dots \to (1) \to \text{head}
\]
For $n < 1$, signal \code{(error 'circ-list "n must be at least 1")}. Use a recursive helper stored in a box, which takes a number \code{k} and a box \code{b}, fills \code{b} with \code{k} and the next box, and builds the rest of the list.
\end{exercise}
\begin{solution}
Two knots are tied here. The head is allocated first as a placeholder, so the last box can point to it. And the helper \code{build} is recursive, so it lives in a box too:
\begin{lstlisting}
(define (circ-list n)
  (if (< n 1)
      (error 'circ-list "n must be at least 1")
      (let* ([head (box 'dummy)]
             [build (box 'dummy)]
             [build-fun
              (lambda (k b)
                (if (equal? k 1)
                    (set-box! b (cons 1 head))
                    (let ([next-box (box 'dummy)])
                      (begin
                        (set-box! b (cons k next-box))
                        ((unbox build) (- k 1) next-box)))))])
        (begin
          (set-box! build build-fun)
          ((unbox build) n head)
          head))))
\end{lstlisting}
\code{build} fills the box it is given, which is the head on the first call, and allocates the next box only when $k > 1$. The box for 1 is filled with \code{head} instead of a new box, which closes the cycle. With the helpers
\begin{lstlisting}
(define (node-value b) (car (unbox b)))
(define (node-next b) (cdr (unbox b)))
\end{lstlisting}
the expected results are
\begin{lstlisting}
(define circ-1 (circ-list 1))
(define circ-3 (circ-list 3))
(check-expect (node-value circ-3) 3)
(check-expect (node-value (node-next circ-3)) 2)
(check-expect (eq? (node-next circ-1) circ-1) #t)
(check-expect (eq? (node-next circ-3) circ-3) #f)
(check-expect (eq? (node-next (node-next (node-next circ-3))) circ-3) #t)
(check-error (circ-list 0) "circ-list: n must be at least 1")
\end{lstlisting}
\code{circ-1} is a one-box cycle, like \code{self-box}, so its next box is itself. Racket prints \code{(circ-list 3)} as \code{\#0='\#\&(3 . \#\&(2 . \#\&(1 . \#0\#)))}.
\end{solution}

\begin{exercise}[Recursive trees]
A tree is \code{(leaf $v$)} or \code{(node $t_1$ $t_2$)}, written as an s-expression \code{'(leaf v)} or \code{'(node left right)}. Write \code{tree-box}, which turns a nested list into a tree: a number becomes a leaf, \code{(list l r)} becomes a node, and anything else is an error \code{"tree: not a tree"}. Tie the knot with a box; \code{tree-box} must not call itself by name.
\begin{center}
\begin{tabular}{@{}ll@{}}
\toprule
input & output \\
\midrule
\code{5} & \code{'(leaf 5)} \\
\code{(list 5 (list 6 7))} & \code{'(node (leaf 5) (node (leaf 6) (leaf 7)))} \\
\code{(list 1 2 3)} & error \\
\bottomrule
\end{tabular}
\end{center}
\end{exercise}
\begin{solution}
The pattern \code{(? number? n)} matches any number and binds it to \code{n}. A node needs two recursive calls, one for each subtree:
\begin{lstlisting}
(define (tree-box input)
  (let* ([tree (box 'dummy)]
         [tree-fun
          (lambda (t)
            (match t
              [(? number? n) (list 'leaf n)]
              [(list left right) (list 'node ((unbox tree) left) ((unbox tree) right))]
              [_ (error 'tree "not a tree")]))])
    (begin
      (set-box! tree tree-fun)
      ((unbox tree) input))))

(check-expect (tree-box 5) '(leaf 5))
(check-expect (tree-box (list 5 (list 6 7)))
              '(node (leaf 5) (node (leaf 6) (leaf 7))))
(check-error (tree-box (list 1 2 3)) "tree: not a tree")
(check-error (tree-box (list 5 (list 6))) "tree: not a tree")
\end{lstlisting}
\code{(list 5 (list 6))} is an error because the inner list has one element, which matches neither a number nor a two-element list. The pitfall above applies here as well: \code{(list 'node ((unbox tree) left) (unbox tree) right)} builds a list of four elements, with the function itself as the third one.
\end{solution}

\subsection{Recursion without Mutation}

Can we get recursion without boxes? The idea is that a function does not refer to itself. Instead, it takes ``the recursive call'' as an argument, as the function passed to \code{fix} in \Cref{sec:cycles} does, but one argument at a time:
\begin{lstlisting}
(define fact-gen
  (lambda (rec)
    (lambda (n)
      (if (equal? 0 n)
          1
          (* n (rec (- n 1)))))))
\end{lstlisting}
\code{fact-gen} is not factorial. Given a function \code{rec} that computes factorial, it returns a function that computes factorial too, by doing one step and leaving the rest to \code{rec}. What we need is a function that is not changed by \code{fact-gen}.

\begin{definition}[Fixed point]
A \emph{fixed point} of a function $f$ is a value $v$ with $f(v) = v$. A \emph{fixed-point combinator} is a function $Y$ such that, for every $f$,
\[
Y\ f = f\ (Y\ f),
\]
that is, $Y\ f$ is a fixed point of $f$.
\end{definition}
If such a $Y$ exists, \code{(Y fact-gen)} is factorial, because unfolding the equation once gives exactly the factorial step:
\begin{lstlisting}[style=repl]
Y fact-gen
=  fact-gen (Y fact-gen)
=  (lambda (n) (if (equal? 0 n) 1 (* n ((Y fact-gen) (- n 1)))))
\end{lstlisting}
Here \code{=} means the two sides are equal, but not by a single evaluation step; \code{=>} below is one step.

\subsubsection{The Y combinator} One such $Y$ is
\begin{lstlisting}
(define Y
  (lambda (h)
    ((lambda (x) (h (x x)))
     (lambda (x) (h (x x))))))
\end{lstlisting}
The trick is \emph{self-application}, \code{(x x)}. The same function is written twice, and the first copy is applied to the second. So inside its body, \code{x} is that function itself, and \code{(x x)} rebuilds the whole application. \code{Y} satisfies the equation:
\begin{lstlisting}[style=repl]
Y f
=> (lambda (x) (f (x x))) (lambda (x) (f (x x)))
=> f ((lambda (x) (f (x x))) (lambda (x) (f (x x))))
=  f (Y f)
\end{lstlisting}
The last step is \code{=}, not \code{=>}: the argument of \code{f} is what \code{(Y f)} stepped to on the first line, but nothing steps back to \code{(Y f)}.

\subsubsection{Why \code{Y} loops in Racket} \code{(define fact-y (Y fact-gen))} never finishes. Write $W$ for \code{(lambda (x) (fact-gen (x x)))}. Then
\begin{lstlisting}[style=repl]
(Y fact-gen)
=> (W W)
=> (fact-gen (W W))
=> (fact-gen (fact-gen (W W)))
=> ...
\end{lstlisting}
Racket evaluates arguments eagerly, so before it can call \code{fact-gen}, it must evaluate the argument \code{(W W)}. That is the expression it started from, one step later. \code{fact-gen} is never called, although it would not even use \code{rec} right away: it would return a \code{lambda} at once.

In a lazy language, \code{Y} works. The argument \code{(W W)} is passed unevaluated, as in the lazy substitution of \Cref{sec:lazy}. \code{fact-gen} returns its \code{(lambda (n) ...)}, and \code{(W W)} is evaluated only when the body calls \code{rec}, which unfolds one more level. A base case that does not call \code{rec} stops the unfolding.

\begin{exercise}[The Z combinator]
Write \code{Z}, the version of \code{Y} for eager evaluation. For every \code{f} and value \code{v} it satisfies
\[
Z\ f\ v = f\ (Z\ f)\ v.
\]
Hint: start from \code{Y}, and delay the self-application by wrapping it in a \code{lambda}: \code{(x x)} becomes \code{(lambda (v) ((x x) v))}.
\end{exercise}
\begin{solution}
This is the $\eta$-expansion of \Cref{sec:recboxes}, applied to both copies of \code{(x x)}:
\begin{lstlisting}
(define Z
  (lambda (h)
    ((lambda (x) (h (lambda (v) ((x x) v))))
     (lambda (x) (h (lambda (v) ((x x) v)))))))

(define fact-z
  (lambda (n) ((Z fact-gen) n)))

(check-expect (fact-z 0) 1)
(check-expect (fact-z 5) 120)
\end{lstlisting}
Write $W'$ for \code{(lambda (x) (fact-gen (lambda (v) ((x x) v))))}. Now
\begin{lstlisting}[style=repl]
(Z fact-gen)
=> (W' W')
=> (fact-gen (lambda (v) ((W' W') v)))
=> (lambda (n) (if (equal? 0 n) 1 (* n (rec (- n 1)))))
       where rec = (lambda (v) ((W' W') v))
\end{lstlisting}
The argument of \code{fact-gen} is a \code{lambda}, which is already a value, so \code{fact-gen} is called, and the result is a function. Evaluation stops there. \code{(W' W')} is evaluated only when \code{rec} is called, and then it produces the same function again, one level down:
\begin{lstlisting}[style=repl]
((Z fact-gen) 3)
=> (* 3 (rec 2))
=> (* 3 ((W' W') 2))
=> (* 3 (* 2 (rec 1)))
=> (* 3 (* 2 (* 1 (rec 0))))
=> (* 3 (* 2 (* 1 1)))
=> 6
\end{lstlisting}
Each line skips the steps that unfold \code{(W' W')} into the factorial function again.
\end{solution}

\begin{exercise}[Putting it together]
Rewrite \code{tree-box} with \code{Z} instead of a box. First write \code{tree-gen}, which takes the recursive call as an argument, like \code{fact-gen}.
\end{exercise}
\begin{solution}
The body of \code{tree-fun} is unchanged, except that the recursive call is the argument \code{rec} instead of \code{(unbox tree)}:
\begin{lstlisting}
(define tree-gen
  (lambda (rec)
    (lambda (t)
      (match t
        [(? number? n) (list 'leaf n)]
        [(list left right) (list 'node (rec left) (rec right))]
        [_ (error 'tree "not a tree")]))))
(define tree-z
  (lambda (t) ((Z tree-gen) t)))
\end{lstlisting}
The four tests of \code{tree-box} pass for \code{tree-z} as well.
\end{solution}

\subsubsection{\code{Z} versus \code{fix}} Both turn a function like \code{fact-gen} into a recursive function. \code{fix} ties the knot with a box: there is one \code{frec}, and it finds itself in the store. \code{Z} uses nothing but functions: the function never finds ``itself'', it rebuilds a copy with \code{(x x)} at every recursive call. (The Racket \code{fix} of \Cref{sec:cycles} takes \code{rec} and \code{n} together, \code{(lambda (rec n) ...)}; \code{Z} takes them one at a time, like the \code{fix} of \Cref{sec:recboxes}.)

\subsection{Challenge: \texttt{rec} via \texttt{Z}}
\label{sec:rec7-interp}

The file \code{recitation7-interp.rkt} contains the interpreter of \Cref{sec:recproblem} with \code{if0C} and \code{recC}, but with no \code{recV} and no boxes: \tyname{Value} has only \code{numV} and \code{closV}. The task is to implement \code{rec} by rewriting it into a call to \code{Z}, as \Cref{sec:recboxes} rewrote it into a call to \code{fix}. A helper \code{run} parses and evaluates:
\begin{lstlisting}
(define (run (s : s-expression)) : Value
  (eval (parse s)))
\end{lstlisting}

\begin{exercise}[Z in our language]
Write \code{Z} as a program in our language. \code{lambda} takes exactly one argument, written \code{(lambda x body)}, and application is \code{(e1 e2)}.
\end{exercise}
\begin{solution}
The Racket \code{Z} already uses one-argument functions only, so only the syntax changes:
\begin{lstlisting}
(define Z
  (parse '(lambda h
            ((lambda x (h (lambda v ((x x) v))))
             (lambda x (h (lambda v ((x x) v))))))))

(define fact-gen
  (parse '(lambda self
            (lambda n
              (if0 n 1 (* n (self (+ n -1))))))))

(test (eval (appC (appC Z fact-gen) (numC 0))) (numV 1))
(test (eval (appC (appC Z fact-gen) (numC 5))) (numV 120))
\end{lstlisting}
In \code{fact-gen}, which the handout provides, the parameter is called \code{self} and not \code{rec}, because \code{rec} is a keyword of our language: \code{(rec (+ n -1))} would parse as a \code{rec} expression.
\end{solution}

\begin{exercise}[\code{rec}]
Fill in the \code{recC} case of \code{eval-env}.
\end{exercise}
\begin{solution}
Exactly as with \code{fix}, \code{(rec f x e)} becomes \code{(Z (lambda f (lambda x e)))}:
\begin{lstlisting}
    [recC (f x e) (eval-env env (appC Z (lambdaC f (lambdaC x e))))]
\end{lstlisting}
\code{(lambda f (lambda x e))} plays the role of \code{fact-gen}. All four tests pass:
\begin{lstlisting}
(test (run '((rec fact x (if0 x 1 (* x (fact (+ x -1))))) 5)) (numV 120))
(test (run '(let fact (rec f x (if0 x 1 (* x (f (+ x -1)))))
              (fact 6)))
      (numV 720))
(test (run '(let fibb (rec f x (if0 x
                                    0
                                    (if0 (+ x -1)
                                         1
                                         (+ (f (+ x -1)) (f (+ x -2))))))
              (fibb 6)))
      (numV 8))
(test (run '(let k 10
              (let f (rec g x (if0 x k (g (+ x -1))))
                (f 3))))
      (numV 10))
\end{lstlisting}
The last test checks that the body of a \code{rec} can still use variables from outside it. It passes for the reason of \Cref{sec:recboxes}: \code{Z} is a closed program, and the user's \code{e} is passed to it as an argument, so \code{Z}'s names \code{h}, \code{x}, and \code{v} never capture anything in \code{e}. That holds even though the tests use \code{x} as their own parameter name.
\end{solution}

\begin{exercise}[Y instead of Z]
Replace \code{Z} by \code{Y} in the \code{recC} case and run the file again. What happens, and why?
\end{exercise}
\begin{solution}
The \code{Z} tests still pass, but the first \code{rec} test never finishes, and memory use grows steadily until DrRacket is stopped. This is the loop of the Racket \code{Y}, one level down. Our language is eager too: \code{appC} evaluates the argument \code{e2} to a value before it runs the body. So \code{(h (x x))} evaluates \code{(x x)}, which evaluates \code{(h (x x))} again, without ever calling \code{h}. Unlike the Racket loop above, this one also grows the Racket stack: each \code{appC} evaluates its argument with a nested call to \code{eval-env}, which is not a tail call.
\end{solution}

\begin{exercise}[Three ways to implement \code{rec}]
For each of \code{recV}, \code{fix}, and \code{Z}: does it need a new kind of \tyname{Value}? Does it need boxes (a store)?
\end{exercise}
\begin{solution}
\begin{center}
\begin{tabular}{@{}>{\raggedright}p{0.2\textwidth}>{\raggedright}p{0.24\textwidth}>{\raggedright}p{0.22\textwidth}>{\raggedright\arraybackslash}p{0.24\textwidth}@{}}
\toprule
& \code{recV} (\Cref{sec:recproblem}) & \code{fix} (\Cref{sec:recboxes}) & \code{Z} (this section) \\
\midrule
new kind of \tyname{Value} & yes, \code{recV} & no & no \\
boxes / a store & no & yes & no \\
interpreter changes & \code{appC} gets a case & only \code{recC} & only \code{recC} \\
where the self-reference is & environment, rebound by each call & a cycle through the store & rebuilt by \code{(x x)} at each call \\
\bottomrule
\end{tabular}
\end{center}
\code{recV} changes the interpreter the most, but it is the most direct: the function's own name is bound to the function, at each call. \code{fix} and \code{Z} change nothing but \code{recC}, because both turn \code{rec} into a program the interpreter can already run. \code{fix} needs the language to have boxes. \code{Z} needs only \code{lambda} and application, which is why recursion needs no special support in an untyped language with first-class functions.
\end{solution}

\begin{remark}
\code{Z} relies on \code{(x x)}, and this self-application cannot be typed directly in a language like \code{plai-typed}: \code{x} would need a type of functions whose argument has that same type. This is the same obstacle as the box that contains itself in \Cref{sec:cycles}. (A recursive \code{define-type} that wraps \code{x} gets around it.) Typed languages therefore usually build recursion in, as \code{define} and \code{letrec} are built into Racket.
\end{remark}
